\documentclass[11pt,a4paper]{article}
\usepackage[utf8]{inputenc}
\usepackage[T1]{fontenc}
\usepackage{amsmath}
\usepackage{booktabs}
\usepackage{listings}
\usepackage[margin=2.3cm]{geometry}

\lstset{
  basicstyle=\ttfamily\footnotesize,
  breaklines=true,
  frame=single,
  numbers=left,
  numberstyle=\tiny,
  showstringspaces=false
}

\begin{document}

% ===================== PAGE DE GARDE =====================
\begin{titlepage}
\centering
\vspace*{1cm}
{\Large \textbf{Institut 2iE}}\\[0.4cm]
{\large Méthodes numériques 2 -- S4/S5D -- 2026--2027}\\[2.5cm]

{\Huge \textbf{TP 01}}\\[0.6cm]
{\LARGE Calculer une énergie avec Scilab}\\[0.4cm]
{\large Méthode des trapèzes et méthode de Simpson}\\[2.5cm]

{\large \textbf{Groupe :}}\\[0.4cm]
\renewcommand{\arraystretch}{1.6}
\begin{tabular}{|c|p{5.5cm}|p{5.5cm}|}
\hline
\textbf{N\textsuperscript{o}} & \textbf{Nom} & \textbf{Prénom}\\
\hline
1 & & \\
\hline
2 & & \\
\hline
3 & & \\
\hline
\end{tabular}

\vfill
{\large Enseignant : Dr Babacar LEYE}\\[0.3cm]
{\large 3 novembre 2026}
\end{titlepage}

% ===================== CONTENU =====================
% ===== PARTIE 1 : objectif et méthodes (Personne A) =====
\section{Objectif}
Une pompe fonctionne pendant 40 minutes et sa puissance a été relevée toutes
les 10 minutes. L'énergie consommée est l'aire sous la courbe de puissance :
\[
  E = \int_0^{T} P(t)\,dt .
\]
Comme on ne dispose que de points mesurés, on estime cette intégrale avec deux
méthodes de quadrature programmées en Scilab : les trapèzes et Simpson.

\section{Méthodes}
\paragraph{Trapèzes.} Sur chaque intervalle $[t_i, t_{i+1}]$, on relie les deux
points par un segment :
\[
  T = \sum_{i=1}^{n} (t_{i+1}-t_i)\,\frac{y_i+y_{i+1}}{2}.
\]
Cette méthode accepte des pas inégaux.

\paragraph{Simpson.} On remplace la courbe par des paraboles sur des paires
d'intervalles :
\[
  S = \frac{h}{3}\Bigl[y_1 + y_{n+1} + 4(y_2+y_4+\dots+y_n)
      + 2(y_3+y_5+\dots+y_{n-1})\Bigr].
\]
Elle exige des pas égaux et un nombre $n$ pair de sous-intervalles. La fonction
programmée refuse les autres cas avec un message d'erreur.

\input{partie2}
\input{partie3}

% ===================== ANNEXE =====================
\appendix
\section{Code Scilab}
\begin{lstlisting}
// Methode des trapezes
function q = trapezes(t, y)
    n = length(t) - 1;
    q = 0;
    for i = 1:n
        h = t(i+1) - t(i);
        if h <= 0 then error("Temps non croissants"); end
        aire = h * (y(i) + y(i+1)) / 2;
        q = q + aire;
    end
endfunction

// Methode de Simpson
function q = simpson(t, y)
    n = length(t) - 1;
    if modulo(n, 2) <> 0 then error("Simpson exige n pair"); end
    h = t(2) - t(1);
    for i = 2:n
        pas = t(i+1) - t(i);
        if abs(pas - h) > 1e-10 then error("Simpson exige des pas egaux"); end
    end
    q = y(1) + y(n+1);
    for i = 2:n
        if modulo(i, 2) == 0 then
            q = q + 4 * y(i);
        else
            q = q + 2 * y(i);
        end
    end
    q = h / 3 * q;
endfunction

// Calcul pour la pompe
tmin = [0 10 20 30 40];
Pkw = [2.0 2.6 3.0 2.5 1.9];
th = zeros(1, length(tmin));
for i = 1:length(tmin)
    th(i) = tmin(i) / 60;
end
ET = trapezes(th, Pkw);
ES = simpson(th, Pkw);
ecart = abs(ET - ES);
\end{lstlisting}

\end{document}
